% Unit 077 terjemahan Bahasa Indonesia dari chapter6.tex baris 1094--1221.
% Sumber: Methods of Algebra, Volume 2, commit
% 9a5803ff77dd3257484cb177f851a73770a59dd3 (CC BY 4.0).
% stable-unit-id: o014.aljabr2.chapter6.group-extensions-revisited
% Cakupan: Bagian 6.6 lengkap.

% segment-id: o014.li2.chapter6.group-extensions-revisited.h001
\section{Meninjau kembali ekstensi grup}\label{sec:group-ext-revisited}
\hypertarget{o014.aljabr2.chapter6.group-extensions-revisited}{ }

% segment-id: o014.li2.chapter6.group-extensions-revisited.p001
Untuk memperoleh pemahaman yang lebih konkret tentang operasi yang
didefinisikan dalam \S\ref{sec:change-of-group}, dan untuk melanjutkan
pembahasan awal dalam \S\ref{sec:grp-low-degree} mengenai ekstensi grup
(Definisi~\ref{def:group-ext}), kita kembali berfokus pada hubungan antara
ekstensi grup dan $\Hm^2$. Kita mulai dari sudut pandang teori grup. Berikut
ini kita tinjau modul-$G$ $M$ dan ekstensi
$0\to M\to E\xrightarrow{\pi}G\to1$. Ekstensi semacam ini mempunyai operasi
tarikan balik, dorongan keluar, dan penjumlahan yang alami, dengan gagasan
serupa dengan \S\ref{sec:Hom-Ext}.

% segment-id: o014.li2.chapter6.group-extensions-revisited.p002
\begin{description}
	\item[Tarikan balik] Misalkan $\varphi:H\to G$ homomorfisme grup. Ambil
	hasil kali serat dalam kategori grup
	\[ \varphi^*E:=E\dtimes{G}H
	=\{(e,h)\in E\times H:\pi(e)=\varphi(h)\}. \]
	Monomorfisme $M\to\varphi^*E$, $x\mapsto(x,1_H)$, menentukan ekstensi grup
	$0\to M\to\varphi^*E\xrightarrow{\text{proyeksi}}H\to1$. Ini memberikan
	funktor $\cate{Ext}(G,M)\to\cate{Ext}(H,M)$. Jika $\varphi$ adalah
	penyertaan subgrup $H$, tarikan balik sepanjang $\varphi$ sama dengan
	mengambil $\pi^{-1}(H)$.

	\item[Dorongan keluar] Misalkan $\theta:M\to N$ homomorfisme modul-$G$.
	Ambil koproduk serat
	\[ \theta_*E:=N\dsqcup{M}E
	=\frac{N\times E}{(y+\theta(x),e)\sim(y,xe)},
	\qquad x\in M,\ y\in N,\ e\in E. \]
	Tuliskan $[y,e]$ untuk citra $(y,e)\in N\times E$ di dalamnya. Perkaliannya
	adalah
	\[ [y_1,e_1][y_2,e_2]
	=[y_1+\pi(e_1)y_2,e_1e_2]. \]
	Terdapat homomorfisme alami
	\[ N\xrightarrow{y\mapsto[y,1_E]}\theta_*E
	\xrightarrow{[y,e]\mapsto\pi(e)}G, \]
	yang menentukan ekstensi grup $\theta_*E$. Ini memberikan funktor
	$\cate{Ext}(G,M)\to\cate{Ext}(G,N)$.

	\item[Mengambil invers] Sebagai kasus khusus, $E$ dapat didorong keluar
	sepanjang homomorfisme modul-$G$ $M\xrightarrow{x\mapsto-x}M$; ekstensi yang
	dihasilkan ditulis sebagai $-E$.

	\item[Jumlah langsung] Untuk objek-objek $E_i$ dalam
	$\cate{Ext}(G_i,M_i)$ untuk $i=1,2$, terdapat ekstensi grup alami
	\[ 0\to M_1\oplus M_2\to E_1\times E_2\to G_1\times G_2\to1, \]
	dengan $G_1\times G_2$ beraksi pada $M_1\oplus M_2$ melalui
	$(g_1,g_2)(x_1,x_2)=(g_1x_1,g_2x_2)$. Tuliskan ekstensi ini sebagai
	$E_1\oplus E_2$.

	\item[Jumlah Baer] Sekarang ambil objek $E_1,E_2$ dalam
	$\cate{Ext}(G,M)$ dan bentuk $E_1\oplus E_2$. Pemetaan penjumlahan
	$M\oplus M\to M$ merupakan homomorfisme modul-$G$ untuk aksi subgrup
	diagonal $G$ dalam $G\times G$. Karena itu, mula-mula tarik balik
	$E_1\oplus E_2$ sepanjang penyertaan diagonal $G\hookrightarrow G\times G$,
	lalu dorong keluar sepanjang homomorfisme penjumlahan $M\oplus M\to M$.
	Objek $\cate{Ext}(G,M)$ yang diperoleh ditulis $E_1\dot{+}E_2$.
	\index{ekstensi grup!jumlah Baer (Baer sum)}
\end{description}

% segment-id: o014.li2.chapter6.group-extensions-revisited.p003
Tarikan balik dan dorongan keluar dapat digabungkan menjadi satu operasi.
Misalkan $\varphi:H\to G$ homomorfisme grup dan homomorfisme grup aditif
$\theta:M\to N$ ekuivarian terhadap $\varphi$
(Definisi~\ref{def:group-equivariance}). Terdapat funktor terkait
$\cate{Ext}(G,M)\to\cate{Ext}(H,N)$, yang didefinisikan dengan terlebih dahulu
menarik balik sepanjang $\varphi$, kemudian mendorong keluar sepanjang
$\theta:\varphi^*M\to N$.

% segment-id: o014.li2.chapter6.group-extensions-revisited.p004
Kompleks standar ternormalisasi $\overline{C}(G,M)$ mempunyai operasi-operasi
sejajar. Untuk $\varphi$ dan $\theta$ seperti di atas,
Proposisi~\ref{prop:grp-equiv} memberikan secara eksplisit morfisme kompleks
$\overline{C}(G,M)\to\overline{C}(H,N)$; dari sini dapat didefinisikan operasi
tarikan balik, dorongan keluar, dan pengambilan invers. Untuk jumlah langsung,
definisinya adalah
\begin{gather*}
	\oplus:\overline{C}^n(G_1,M_1)\oplus\overline{C}^n(G_2,M_2)
	\to\overline{C}^n(G_1\times G_2,M_1\oplus M_2),\\
	f_1\oplus f_2:((g_{1,i},g_{2,i}))_{i=1}^n\mapsto
	\left(f_1(g_{1,1},\ldots,g_{1,n}),
	f_2(g_{2,1},\ldots,g_{2,n})\right).
\end{gather*}
Ketika $n$ bervariasi, ini memberikan morfisme kompleks
$\overline{C}(G_1,M_1)\oplus\overline{C}(G_2,M_2)
\to\overline{C}(G_1\times G_2,M_1\oplus M_2)$. Dengan pola yang sama,
definisikan jumlah Baer pada $\overline{C}(G,M)$ memakai jumlah langsung,
tarikan balik, dan dorongan keluar. Hasilnya tepat merupakan penjumlahan
dalam $\overline{C}^n(G,M)$.

% segment-id: o014.li2.chapter6.group-extensions-revisited.t001
\begin{proposition}\label{prop:factor-set-operation}
	Hingga isomorfisme kanonik, operasi tarikan balik, dorongan keluar,
	pengambilan invers, jumlah langsung, dan jumlah Baer pada ekstensi grup
	bersesuaian, di bawah ekuivalensi Teorema~\ref{prop:factor-set}, dengan
	operasi-operasi terkait pada kompleks
	$\tau^{\leq2}\overline{C}(G,M)$ (atau kategori
	$\tau^{\leq2}\overline{\cate{C}}(G,M)$).
\end{proposition}
\begin{proof}
	Untuk suatu ekstensi grup $0\to M\to E\to G\to1$, pilih penampang
	ternormalisasi $s:G\to E$ dan $f\in\overline{Z}^2(G,M)$ terkait seperti
	dalam bukti Teorema~\ref{prop:factor-set}. Misalkan $\varphi:H\to G$
	homomorfisme grup. Penampang ternormalisasi untuk $\varphi^*E$ dapat diambil
	sebagai
	\[ s':h\mapsto(s(\varphi(h)),h)\in E\dtimes{G}H. \]
	Maka $s'(h_1)s'(h_2)=f'(h_1,h_2)s'(h_1h_2)$, dengan
	\[ f'(h_1,h_2):=f(\varphi(h_1),\varphi(h_2)),\qquad h_1,h_2\in H. \]
	Ini membuktikan kasus tarikan balik. Untuk dorongan keluar sepanjang
	homomorfisme modul-$G$ $\theta:M\to N$, dengan notasi di atas ambil
	penampang ternormalisasi dari $\theta_*E$ sebagai
	\[ s'(g):=[0,s(g)]\in\theta_*E=N\dsqcup{M}E. \]
	Perhitungan dalam $\theta_*E$ memberikan
	\begin{align*}
		s'(g_1)s'(g_2)&=[0,s(g_1)][0,s(g_2)]=[0,s(g_1)s(g_2)]\\
		&=[0,f(g_1,g_2)s(g_1g_2)]
		=[\theta f(g_1,g_2),s(g_1g_2)]\\
		&=[\theta f(g_1,g_2),1_E]s'(g_1g_2).
	\end{align*}
	Jadi $f'$ yang bersesuaian dengan $s'$ tepat merupakan $\theta f$.
	Ini membuktikan kasus dorongan keluar.

	Kasus jumlah langsung diperiksa dengan cara yang sama dan
	perinciannya diserahkan kepada pembaca. Pengambilan invers dan jumlah Baer
	dapat diuraikan menjadi operasi-operasi di atas.
\end{proof}

% segment-id: o014.li2.chapter6.group-extensions-revisited.p005
Sebagai kasus khusus, hingga isomorfisme kanonik, pengambilan invers dan jumlah
Baer pada ekstensi grup masing-masing bersesuaian dengan invers dan
penjumlahan dalam grup $\Hm^2(G,M)$.

% segment-id: o014.li2.chapter6.group-extensions-revisited.d001
\begin{definition}\label{def:central-ext}
	\index{ekstensi grup!pusat (central)}
	Jika dalam ekstensi grup $0\to A\to E\to G\to1$, grup $A$ merupakan
	subgrup pusat dari $E$, maka ekstensi tersebut disebut \emph{ekstensi
	pusat}.

	Untuk suatu ekstensi pusat seperti di atas, jika untuk setiap ekstensi pusat
	$0\to B\to F\to G\to1$ terdapat homomorfisme grup $\varphi:E\to F$ tunggal
	yang membuat diagram berikut komutatif,
	\[\begin{tikzcd}
		0 \arrow[r] & A \arrow[r] \arrow[d,"{\varphi|_A}"'] &
		E \arrow[d,"\varphi"] \arrow[r] & G \arrow[equal,d] \arrow[r] & 1 \\
		0 \arrow[r] & B \arrow[r] & F \arrow[r] & G \arrow[r] & 1,
	\end{tikzcd}\]
	maka $0\to A\to E\to G\to1$ disebut \emph{ekstensi pusat universal} dari
	$G$. Jika ekstensi ini ada, ia unik hingga isomorfisme tunggal.
\end{definition}

% segment-id: o014.li2.chapter6.group-extensions-revisited.p006
Menurut definisi, semua automorfisme dalam dari ekstensi pusat bersifat
trivial. Untuk setiap grup aditif $A$, definisikan kategori ekstensi pusat
terkait
\[ \cate{CExt}(G,A):=\cate{Ext}(G,A),\qquad
	\text{dengan $A$ dipandang sebagai modul-$G$ beraksi trivial}. \]

% segment-id: o014.li2.chapter6.group-extensions-revisited.l001
\begin{lemma}\label{prop:univ-central-ext-prep}
	Jika grup $G$ mempunyai ekstensi pusat universal
	$0\to A\to E\xrightarrow{\pi}G\to1$, maka
	$G=G_{\mathrm{der}}$ dan $E=E_{\mathrm{der}}$.
\end{lemma}
\begin{proof}
	Perhatikan ekstensi pusat terbelah
	$0\to E_{\mathrm{ab}}\to E_{\mathrm{ab}}\times G\to G\to1$ dari $G$.
	Terdapat homomorfisme $(0,\pi)$ dan $(q,\pi)$ dari ekstensi $E$ ke
	$E_{\mathrm{ab}}\times G$, dengan $q:E\to E_{\mathrm{ab}}$ homomorfisme
	hasil bagi. Bagian ketunggalan dari sifat universal mengakibatkan $q=0$,
	yakni $E=E_{\mathrm{der}}$. Dari sini segera diperoleh
	$G=G_{\mathrm{der}}$.
\end{proof}

% segment-id: o014.li2.chapter6.group-extensions-revisited.p007
Berdasarkan konstruksi dorongan keluar, tidak sulit membuktikan bahwa
$\varphi$ dalam Definisi~\ref{def:central-ext} berfaktor secara tunggal
melalui $E\to(\varphi|_A)_*E$. Dengan demikian sifat universal dapat ditulis
ulang sebagai berikut: untuk setiap ekstensi pusat
$0\to B\to F\to G\to1$, terdapat homomorfisme grup $\theta:A\to B$ tunggal
dan isomorfisme $\theta_*E\rightiso F$ tunggal dalam $\cate{CExt}(G,B)$.
Dengan syarat $G=G_{\mathrm{der}}$, kategori $\cate{CExt}(G,B)$ tidak
mempunyai automorfisme taktrivial. Definisi ekstensi pusat universal dapat
ditulis lagi: untuk setiap grup abelian $B$, terdapat bijeksi
\[\begin{tikzcd}[row sep=tiny]
	\Hom_{\cate{Ab}}(A,B) \arrow[r,"1:1"] &
	\Obj(\cate{CExt}(G,B))\big/\simeq \\
	\theta \arrow[mapsto,r] & \text{kelas isomorfisme dari }\theta_*E.
\end{tikzcd}\]
Ruas kiri merupakan funktor dalam $B$. Ruas kanan juga memberikan funktor
$\kappa:\cate{Ab}\to\cate{Set}$, dengan kefunktoran dalam $B$ diberikan oleh
dorongan keluar ekstensi. Jadi keberadaan ekstensi pusat universal setara
dengan pertanyaan apakah $\kappa$ dapat direpresentasikan.

% segment-id: o014.li2.chapter6.group-extensions-revisited.t002
\begin{theorem}
	Grup $G$ mempunyai ekstensi pusat universal jika dan hanya jika
	$G=G_{\mathrm{der}}$. Jika syarat ini berlaku, subgrup pusat $A$ dalam
	ekstensi pusat universal isomorfik dengan $\Hm_2(G,\Z)$.
\end{theorem}
\begin{proof}
	Menurut Lema~\ref{prop:univ-central-ext-prep}, cukup buktikan arah
	``jika''. Perhatikan funktor $\kappa$ di atas. Teorema~\ref{prop:factor-set}
	mengidentifikasi $\kappa(B)$ dengan $\Hm^2(G,B)$, sedangkan
	Proposisi~\ref{prop:factor-set-operation} menunjukkan bahwa dorongan keluar
	bersesuaian dengan operasi terkait pada $\Hm^2(G,\cdot)$. Kita mengetahui
	$\Hm_1(G,\Z)\simeq G_{\mathrm{ab}}=0$
	(Contoh~\ref{eg:group-H1-further}). Jadi dengan mengambil $n=2$ dalam
	Proposisi~\ref{prop:group-UCT}, diperoleh
	$\Hm^2(G,B)\simeq\Hom_{\cate{Ab}}(\Hm_2(G,\Z),B)$, dan isomorfisme ini
	natural dalam $B$. Dengan demikian $\kappa$ dapat direpresentasikan.
\end{proof}

% segment-id: o014.li2.chapter6.group-extensions-revisited.p008
Grup $\Hm_2(G,\Z)$ juga disebut \emph{pengali Schur} dari $G$.
Korolari~\ref{prop:Hopf-H2} di bawah berguna untuk menghitungnya. Dalam praktik,
hal terpenting sering kali adalah deskripsi eksplisit ekstensi pusat universal;
karakterisasi keberadaan di atas baru merupakan langkah pertama.
\index{pengali Schur (Schur multiplier)}

% segment-id: o014.li2.chapter6.group-extensions-revisited.p009
Sebagai penerapan lain dari $\Hm^2$, kita turunkan sebuah hasil teori grup.

% segment-id: o014.li2.chapter6.group-extensions-revisited.t003
\begin{theorem}[Schur--Zassenhaus]\label{prop:Schur-Zassenhaus}
	Misalkan diberikan ekstensi grup $1\to A\to E\xrightarrow{\pi}G\to1$, dengan $A$
	dan $G$ grup hingga yang ordenya relatif prima. Maka ekstensi ini terbelah.
\end{theorem}
\begin{proof}
	Langkah pertama adalah menangani kasus ketika $A$ abelian. Berikan $A$
	struktur modul-$G$ yang berasal dari ekstensi grup. Cukup dibuktikan bahwa
	$\Hm^2(G,A)$ trivial. Karena $|A|$ dan $|G|$
	relatif prima, cukup buktikan kesamaan berikut untuk $k=2$:
	\[ k\geq1\implies |A|\cdot\Hm^k(G,A)=0
	=|G|\cdot\Hm^k(G,A). \]

	Kesamaan pertama mudah. Untuk setiap $t\in\Z$, endomorfisme $A\to A$ yang
	diberikan oleh perkalian dengan $t$ menginduksi endomorfisme pada setiap
	$\Hm^k(G,A)$ yang juga berupa perkalian dengan $t$; ini merupakan latihan
	sederhana dalam bab ini. Kesamaan kedua juga merupakan latihan dalam bab
	ini; Korolari~\ref{prop:group-coh-torsion} di bawah akan memberikan bukti yang
	lengkap, meskipun agak memutar.

	Untuk $A$ umum, tujuan berikutnya adalah menunjukkan bahwa $E$ mempunyai
	subgrup $H$ berorde $|G|$. Ini akan memberikan $H\cap A=\{1\}$ dan
	$\pi|_H:H\rightiso G$; inversnya memberi pembelahan ekstensi grup. Kita
	menggunakan induksi pada $|E|$. Tanpa mengurangi keumuman, anggap $|A|>1$.

	Ambil bilangan prima $p$ yang membagi $|A|$, sehingga $p\nmid|G|$, lalu
	ambillah subgrup Sylow-$p$ $P$ dari $A$. Maka $P$ juga merupakan subgrup
	Sylow-$p$ dari $E$. Karena $A\lhd E$ dan semua subgrup Sylow-$p$ saling
	konjugat, semuanya terkandung dalam $A$. Dengan menghitung banyaknya
	subgrup Sylow-$p$ masing-masing dalam $A$ dan $E$, diperoleh
	$(E:N_E(P))=(A:N_A(P))$, dengan $N_E$ dan $N_A$ menyatakan penormal; atau,
	setelah disusun ulang,
	\[ (N_E(P):N_A(P))=(E:A)=|G|. \]

	Berlaku $N_A(P)\lhd N_E(P)$. Ekstensi grup terkait
	\[ 1\to N_A(P)\to N_E(P)\to N_E(P)/N_A(P)\to1 \]
	masih memenuhi syarat teorema. Jika $N_E(P)\neq E$, hipotesis induksi
	memberikan subgrup $H$ dari $N_E(P)$ dengan
	$|H|=(N_E(P):N_A(P))=|G|$, sehingga pernyataan terbukti.

	Sekarang kita boleh menganggap $P\lhd E$, sehingga $P\lhd A$. Perhatikan
	ekstensi grup
	\[ 1\to A/P\to E/P\to G\to1. \]
	Menurut hipotesis induksi, terdapat subgrup $H$ dari $E$ dengan
	$H\supset P$ dan $|H/P|=|G|$. Fakta dasar teori grup menunjukkan bahwa
	pusat $Z:=Z_P$ dari $P$ taktrivial. Kita mempunyai $Z\lhd H$ dan
	$(H/Z:P/Z)=(H:P)=|G|$. Dengan kembali memakai induksi, terdapat subgrup
	$K$ dari $H$ dengan $K\supset Z$ dan $|K/Z|=|G|$.

	Sekarang terdapat ekstensi grup $1\to Z\to K\to K/Z\to1$. Karena $Z$
	grup-$p$, jika $K\neq E$ maka induksi memberikan subgrup dari $K$ berorde
	$|K/Z|=|G|$, dan pernyataan terbukti. Selanjutnya anggap $K=E$. Maka
	harus berlaku $H=E$, sehingga dari $(H:P)=|G|$ diperoleh $P=A$; dalam
	keadaan ini $Z=Z_A\lhd E$.

	Terakhir, perhatikan ekstensi grup
	$1\to A/Z\to E/Z\to G\to1$. Kita memperoleh subgrup $Q$ dari $E$ dengan
	$Q\supset Z$ dan $|Q/Z|=|G|$. Jika $A$ tidak abelian, maka $Z\neq A$ dan
	$Q\neq E$; dengan meninjau ekstensi grup
	$1\to Z\to Q\to Q/Z\to1$, kita mendapatkan subgrup dari $Q$ berorde $|G|$.
	Jika $A$ abelian, masalahnya sudah diselesaikan pada awal bukti.
\end{proof}

% segment-id: o014.li2.chapter6.group-extensions-revisited.p010
Jika $A$ abelian, semua pembelahan ekstensi $E$ dalam
Teorema~\ref{prop:Schur-Zassenhaus} saling konjugat oleh $A$. Memang,
menentukan pembelahan setara dengan menentukan isomorfisme ekstensi grup
$A\rtimes G\simeq E$, sehingga dua pembelahan $s,s':G\to E$ berbeda melalui
suatu automorfisme $E$. Akan tetapi, Lema~\ref{prop:factor-set-prep} atau
Teorema~\ref{prop:factor-set} menunjukkan bahwa grup automorfisme luar
ekstensi tersebut adalah $\Hm^1(G,A)$, dan langkah pertama bukti
Teorema~\ref{prop:Schur-Zassenhaus} telah menunjukkan bahwa grup ini trivial.

% segment-id: o014.li2.chapter6.group-extensions-revisited.p011
Sifat konjugasi pembelahan dapat diperluas ke $A$ umum, tetapi pembuktiannya
lebih berliku.
